\chapter{Proposed Scheme}\label{ch:design}

This chapter describes the scheme as proposed at the Phase~I review and as implemented, called
\emph{Phase-I} in the evaluation. Chapter~\ref{ch:opt} adds the optimisations.

\section{Architecture and Design Principle}

\begin{figure}[htbp]\centering
\begin{tikzpicture}[
  box/.style={draw, rounded corners, align=center, minimum width=3.2cm, minimum height=1.25cm, font=\small},
  arr/.style={-{Latex[length=2mm]}, thick}, lab/.style={font=\scriptsize, align=center}]
  \node[box, fill=blue!8] (ta) at (0,4.2) {Trusted Authority\\\scriptsize registration, escrow};
  \node[box, fill=green!8] (ledger) at (7.5,4.2) {Permissioned blockchain\\\scriptsize Hyperledger Fabric};
  \node[box, fill=orange!10] (rsu1) at (3.2,1.6) {RSU $A$\\\scriptsize verify proof, policy};
  \node[box, fill=orange!10] (rsu2) at (8.6,1.6) {RSU $B$\\\scriptsize handover, detect};
  \node[box, fill=gray!10] (v1) at (1.5,-1.2) {Vehicle (OBU)\\\scriptsize pseudonym $\mathit{pid}_i$};
  \node[box, fill=gray!10] (v2) at (6.8,-1.2) {Vehicle (OBU)};
  \draw[arr] (ta) -- node[lab, above]{Merkle root,\\policies} (ledger);
  \draw[arr, dashed] (ta) -- node[lab, left]{pseudonyms,\\keys (offline)} (v1);
  \draw[arr] (v1) -- node[lab, left]{ZKP login} (rsu1);
  \draw[arr] (rsu1) -- node[lab, right, pos=0.35]{ticket} (v1);
  \draw[arr, <->] (rsu1) -- node[lab, left, pos=0.55]{root, policy, auth,\\reports, revocations} (ledger);
  \draw[arr, <->] (rsu2) -- (ledger);
  \draw[arr, <->] (rsu1) -- node[lab, above]{PREAUTH, QUARANTINE\\(wired backbone)} (rsu2);
  \draw[arr] (v1) -- node[lab, below]{signed BSMs (10\,Hz)} (v2);
  \draw[arr] (v2) -- (rsu2);
\end{tikzpicture}
\caption{Entities and the information each exchanges. The blockchain holds only state that several
RSUs must share and that must be auditable; it is not on the path of individual safety messages.}
\label{fig:architecture}
\end{figure}

Figure~\ref{fig:architecture} shows the entities. The guiding principle of the design is that
\emph{proofs are verified at the edge and the ledger stores only shared state}. An RSU verifies a
zero-knowledge proof locally in well under a millisecond, while a Fabric commit takes hundreds of
milliseconds (Section~\ref{sec:res-ledger}). So the ledger is read from a local, regularly
synchronised copy, written asynchronously, and never waited for while a safety message is handled.
The ledger stores:
\begin{itemize}
  \item the Merkle root of each registration epoch (never a per-vehicle record);
  \item service policies: which vehicle classes may use which service;
  \item authorization records keyed by pseudonym, so that the next RSU can accept a handover;
  \item misbehaviour reports with evidence hashes, and the revocation list;
  \item trace requests and the RSU organizations' approvals for identity reveal.
\end{itemize}

\section{Registration}

For each vehicle $v$ the TA generates $P$ pseudonyms (5 by default), each a fresh key pair
$(x_i, Y_i = x_i G)$ with a random 64-bit identifier $\mathit{pid}_i$, and the vehicle's class
(\textit{car}, \textit{truck} or \textit{emergency}). It builds the Merkle tree over the leaves of
all pseudonyms of all vehicles, in shuffled order so that positions do not reveal which pseudonyms
belong together, and gives each vehicle its keys and Merkle paths over a secure channel. It publishes
the root with its validity period on the ledger through the chaincode function
\texttt{RegisterRoot}. Finally it encrypts the vehicle's identity record (identity and list of
pseudonyms) and escrows the key with the RSU organizations as $t$-of-$n$ Shamir shares, with
$t = 2f+1$ for $n = 3f+1$ organizations (5 of 9 in the evaluated scenarios).

\section{Authentication}

\begin{figure}[htbp]\centering
\begin{tikzpicture}[font=\small, >={Latex[length=2mm]}]
  \node (v) at (0,0) {\textbf{Vehicle}};
  \node (r) at (6,0) {\textbf{RSU}};
  \node (l) at (11,0) {\textbf{Ledger}};
  \draw[gray] (0,-0.3) -- (0,-8.3); \draw[gray] (6,-0.3) -- (6,-8.3); \draw[gray] (11,-0.3) -- (11,-8.3);
  \draw[->] (6,-0.8) -- node[above, font=\scriptsize]{BEACON: $\mathit{rsu},\ n,\ ts,\ \sigma_{RSU}$ (every 100 ms)} (0,-0.8);
  \node[align=left, font=\scriptsize, anchor=west] at (0.1,-1.6) {verify beacon; $c = H(G\|Y\|T\|\mathit{ctx})$,\\ $s = r + cx$ with a precomputed $T = rG$};
  \draw[->] (0,-2.4) -- node[above, font=\scriptsize]{AUTH\_REQ: $\mathit{pid}, Y, \mathit{path}, (T,s), \mathit{svc}, \mathit{class}, n, ts$} (6,-2.4);
  \node[align=left, font=\scriptsize, anchor=west] at (6.1,-3.35) {nonce fresh, unused\\ Merkle path $\to$ root\\ $sG = T + cY$\\ not revoked, policy allows};
  \draw[dashed, ->] (11,-3.0) -- node[above, font=\scriptsize]{root, revocations, policy (local copy)} (6,-3.0);
  \draw[->] (6,-4.6) -- node[above, font=\scriptsize]{AUTH\_RESP: ticket $\{\mathit{pid}, Y, \mathit{svc}, \mathit{class}, \mathit{exp}\}_{\sigma_{RSU}}$} (0,-4.6);
  \draw[->] (6,-5.4) -- node[above, font=\scriptsize]{RecordAuth (asynchronous)} (11,-5.4);
  \draw[->] (6,-6.1) -- node[above, font=\scriptsize]{PREAUTH to neighbouring RSUs (backbone)} (9,-6.1);
  \draw[->] (0,-7.0) -- node[above, font=\scriptsize]{BSM: $\mathit{pid}$, kinematics, ticket (or its digest), $\mathrm{Sig}_{x}$} (6,-7.0);
  \node[align=left, font=\scriptsize, anchor=west] at (6.1,-7.7) {ticket signature, pid match,\\ BSM signature, detector};
\end{tikzpicture}
\caption{Authentication, ticket issue and the first safety message}\label{fig:auth-seq}
\end{figure}

Figure~\ref{fig:auth-seq} shows the exchange.
\begin{enumerate}
  \item Every RSU broadcasts a signed beacon every 100\,ms with its identity, a fresh nonce $n$ and a
  timestamp. A vehicle accepts only beacons from RSUs whose public keys it holds.
  \item The vehicle computes the Schnorr response for the current pseudonym with
  \[\mathit{ctx} = \mathsf{enc}(\texttt{"VANET-AUTH-v1"}, \mathit{rsu}, n, ts, \mathit{svc}, \mathit{pid}, \mathit{class})\]
  and sends the pseudonym, its public key, the Merkle path, the proof, the requested service and its
  class.
  \item The RSU checks that the nonce is one it issued recently and that this (nonce, pseudonym) pair
  was not used before, rebuilds the leaf and verifies the path against the current root from the
  ledger, verifies the proof, checks that the pseudonym is neither revoked nor quarantined, and
  evaluates the authorization predicate
  \begin{equation}\label{eq:access}
    \mathrm{Access}(v, \mathit{svc}) = 1 \iff \mathrm{Verify}(\pi, Y) = 1 \;\wedge\; \mathit{pid} \notin \mathit{RL}
      \;\wedge\; \mathit{class} \in \mathrm{Policy}(\mathit{svc}).
  \end{equation}
  \item If access is granted, the RSU signs a ticket over the pseudonym, its key, the service, the
  class, its own identity and an expiry 30\,s ahead, records the authorization on the ledger
  asynchronously, and sends a signed PREAUTH note over the backbone to its neighbouring RSUs.
\end{enumerate}
The RSU learns the pseudonym, the class and that the vehicle is registered, and nothing else: the
leaf of another vehicle of the same class looks the same, the proof is zero-knowledge, and the
Merkle path reveals only the leaf position, which is shuffled.

\section{Handover and Pseudonym Change}

When a vehicle approaches the edge of its RSU's range (80\,\% of the nominal 300\,m) and hears a
closer RSU, or loses its RSU, it sends a handover request with its ticket, signed with the ticket key
and bound to the new RSU's beacon nonce. The new RSU checks the issuing RSU's signature on the ticket,
the request signature and the shared authorization: either the authorization record on the ledger or,
if the commit has not reached its local copy yet, the PREAUTH note from the issuing RSU. A handover
needs no new proof. After three sessions a vehicle switches to its next pseudonym and authenticates
afresh, so that a pseudonym's lifetime stays short.

\section{Safety Messages}

Each BSM carries the pseudonym, the kinematic fields and the ticket, and is signed with the
pseudonym's secret key, whose public key is in the ticket. A receiver checks the ticket's RSU
signature, that the ticket's pseudonym matches, the BSM signature and the ticket expiry. The receiver
needs no certificate chain and no ledger access. In Phase~I every BSM carries the full ticket, which
makes it about 330 bytes.

\section{Misbehaviour Detection and Reporting}\label{sec:detect}

Every RSU runs a detector on each authentic BSM against that pseudonym's recent history. Hard
plausibility rules flag positions beyond radio range, speeds above 60\,m/s, accelerations above the
physical limit, positions inconsistent with the claimed speed, impossible heading flips, replays and
message rates above 13 per second. Weak rules (sudden appearance close to the RSU, two pseudonyms at
the same place) count only in pairs. An RSU reports a pseudonym once $K = 3$ of its last $W = 10$
messages were flagged, with the hash of a flagged, \emph{signed} message as evidence, through the
chaincode function \texttt{ReportMisbehaviour}.

\section{Revocation and Threshold Identity Reveal}

The chaincode revokes a pseudonym automatically once two distinct RSUs have reported it. A single
report, perhaps a detector error, therefore has no effect. RSUs see the revocation through their
local ledger copy and reject the pseudonym's BSMs and logins. An RSU that reported a now-revoked
pseudonym opens a trace request on the ledger (\texttt{RequestTrace}), and the RSU organizations
approve it (\texttt{ApproveTrace}). Once $t$ approvals are recorded, each approving organization
releases its key share, the TA reconstructs the escrow key, decrypts the identity record and revokes
all of the vehicle's pseudonyms. The vehicle's next login attempt fails and it is shut out. No single
party can unmask a vehicle: the TA lacks the key, and fewer than $t$ RSU organizations hold too few
shares.

\section{Ledger Interface}

\begin{table}[htbp]\centering\small
\caption{Chaincode functions and who may call them}\label{tab:chaincode}
\begin{tabular}{@{}llp{7.2cm}@{}}
\toprule
Function & Caller & Effect \\
\midrule
\texttt{RegisterRoot} & TA & publish the Merkle root of an epoch \\
\texttt{SetPolicy}, \texttt{SetConfig} & TA & classes allowed per service; $K$, trace threshold \\
\texttt{RecordAuth} & RSU & authorization record of a pseudonym (rejected if revoked) \\
\texttt{ReportMisbehaviour} & RSU & report with evidence hash; revoke at $K$ reporters \\
\texttt{Revoke} & TA & revoke pseudonyms after a trace \\
\texttt{RequestTrace}, \texttt{ApproveTrace} & RSU & open and approve an identity reveal \\
\texttt{SetReporterWeight} & TA & report weight of an RSU (Section~\ref{sec:trust}) \\
\texttt{IsRevoked}, \texttt{GetRoot}, \ldots & any & queries; \texttt{ListRevokedSince} for synchronisation \\
\bottomrule
\end{tabular}
\end{table}

Table~\ref{tab:chaincode} lists the chaincode functions. The chaincode reads the caller's role from
its Fabric certificate, so an RSU cannot change a policy and a vehicle cannot write at all. RSUs reach
Fabric through a REST gateway that signs each transaction with the caller's identity.

\section{Security Analysis}

\begin{itemize}
  \item \textbf{Impersonation.} An accepting proof requires $x_i$ for a leaf in the registered tree
  (soundness of the Schnorr proof and collision resistance of the hash). A signed BSM requires the
  key in an RSU-signed ticket.
  \item \textbf{Replay.} A login proof is bound to the RSU, a nonce the RSU issued within the
  freshness window and the pseudonym, and each (nonce, pseudonym) pair is accepted once. BSMs carry a
  timestamp and sequence number that the detector checks; handover requests are bound to the new
  RSU's nonce.
  \item \textbf{Class and service escalation.} The class is part of the Merkle leaf and of the proof
  context, so claiming another class breaks verification; the policy check (\ref{eq:access}) then
  decides the service.
  \item \textbf{Anonymity.} RSUs and vehicles see pseudonyms only, and the proof is zero-knowledge.
  Linking pseudonyms of one vehicle requires its escrow record, which needs $t$ organizations.
  \item \textbf{Accountability.} Every revocation and every reveal is a ledger transaction signed by
  identified organizations, and the ledger is append-only.
  \item \textbf{Limitations.} The TA knows at registration which pseudonyms belong to which vehicle.
  The RSUs learn which pseudonym was where, and a global eavesdropper can link consecutive pseudonyms
  by position. The privacy profile of Chapter~\ref{ch:opt} addresses both.
\end{itemize}
